\documentclass[../../hippo]{subfiles}
\begin{document}

\subsection{Derivation for Translated Chebyshev}
\label{sec:derivation-chebyshev}

The final family of orthogonal polynomials we analyze under the HiPPO framework are the Chebyshev polynomials.
The Chebyshev polynomials can be seen as the purely real analog of the Fourier basis;
for example, a Chebyshev series is related to a Fourier cosine series through a change of basis~\citep{boyd2001chebyshev}.


\paragraph{Measure and Basis}


The basic Chebyshev measure is $\omega^{\mathrm{cheb}} = (1-x^2)^{-1/2}$ on $(-1, 1)$.
Following \cref{sec:chebyshev-properties}, we choose the following measure and orthonormal basis polynomials in terms of the Chebyshev polynomials of the first kind $T_n$.
\begin{align*}%
  \omega(t, x) &= \frac{2}{\theta\pi} \omega^{cheb}\left( \frac{2(x-t)}{\theta} + 1 \right) \mathbb{I}_{(t-\theta, t)}
  \\
  &= \frac{1}{\theta\pi} \left( \frac{x-t}{\theta} + 1 \right)^{-1/2} \left( -\frac{x-t}{\theta} \right)^{-1/2} \mathbb{I}_{(t-\theta, t)}
  \\
  p_n(t, x) &= \sqrt{2} T_n\left( \frac{2(x-t)}{\theta} + 1 \right) \qquad \text{for } n \geq 1,
  \\
  p_0(t, x) &= T_0\left( \frac{2(x-t)}{\theta} + 1 \right).
\end{align*}

Note that at the endpoints, these evaluate to
\begin{align*}%
  p_n(t, t) &=
  \begin{cases}
    \sqrt{2} T_n(1) = \sqrt{2} & n \ge 1
    \\
    T_n(1) = 1 & n = 0
  \end{cases}
  \\
  p_n(t, t-\theta) &=
  \begin{cases}%
    \sqrt{2} T_n(-1) = \sqrt{2} (-1)^{n} & n \ge 1
    \\
    T_n(-1) = 1 & n = 0
  \end{cases}
\end{align*}

\paragraph{Tilted Measure}


Now we choose
\begin{align*}
  \chi^{(t)} = 8^{-1/2} \theta \pi \omega^{(t)},
\end{align*}
So
\begin{align*}
  \frac{\omega}{\chi^2} = \frac{1}{\frac{\theta^2\pi^2}{8}\omega} = \frac{8}{\theta \pi} \left( \frac{x-t}{\theta} + 1 \right)^{1/2} \left( -\frac{x-t}{\theta} \right)^{1/2} \mathbb{I}_{(t-\theta, t)}
\end{align*}
which integrates to $1$.


We also choose $\lambda_n = 1$ for the canonical orthonormal basis, so
\begin{align*}
  g^{(t)} = p_n^{(t)} \chi^{(t)}
\end{align*}


\paragraph{Derivatives}

The derivative of the density is
\begin{align*}
  \ppt \frac{\omega}{\chi} = \ppt \frac{8^{1/2}}{\theta\pi}\mathbb{I}_{(t-\theta, t)} = \frac{8^{1/2}}{\theta\pi}(\delta_t - \delta_{t-\theta}).
\end{align*}

We consider differentiating the polynomials separately for $n = 0$, $n$ even, and $n$ odd,
using equation~\eqref{eq:chebyshev-derivative}.
Defined $z = \frac{2(x-t)}{\theta} + 1$ for convenience.
First, for $n$ even,
\begin{align*}%
  \ppt p_n(t, x) &= -\frac{2^{\frac{3}{2}}}{\theta} T_n'\left( \frac{2(x-t)}{\theta} + 1 \right)
  \\
  &= -\frac{2^{\frac{3}{2}}}{\theta} T_n'\left( z \right)
  \\
  &= -\frac{2^{\frac{3}{2}}}{\theta} \cdot 2n \left( T_{n-1}(z) + T_{n-3}(z) + \dots + T_1(z) \right)
  \\&
  = -\frac{4n}{\theta} \left( p_{n-1}(t, x) + p_{n-3}(t, x) + \dots + p_1(t, x) \right)
\end{align*}
For $n$ odd,
\begin{align*}%
  \ppt p_n(t, x) &= -\frac{2^{\frac{3}{2}}}{\theta} T_n'\left( \frac{2(x-t)}{\theta} + 1 \right)
  \\
  &= -\frac{2^{\frac{3}{2}}}{\theta} T_n'\left( z \right)
  \\
  &= -\frac{2^{\frac{3}{2}}}{\theta} \cdot 2n \left( T_{n-1}(z) + T_{n-3}(z) + \dots + T_1(z) + \frac{1}{2}T_0(z) \right)
  \\&
  = -\frac{4n}{\theta} \left( p_{n-1}(t, x) + p_{n-3}(t, x) + \dots + 2^{-\frac{1}{2}}p_0(t, x) \right)
\end{align*}
And
\begin{align*}
  \ppt p_0(t, x) &= 0.
\end{align*}



\paragraph{Coefficient Dynamics}

\begin{align*}%
  c_n(t) &= \int f(x) p_n(t, x) \frac{2^{3/2}}{\theta\pi} \mathbb{I}_{(t-\theta, t)} \dd x
  \\
  \ddt c_n(t) &= \int f(x) \ppt p_n(t, x) \frac{2^{3/2}}{\theta\pi} \mathbb{I}_{(t-\theta, t)} \dd x + \frac{2^{3/2}}{\theta\pi} f(t)p_n(t, t) - \frac{2^{3/2}}{\theta\pi} f(t-\theta)p_n(t, t-\theta)
  \\&
  = -\frac{4n}{\theta} (c_{n-1} + c_{n-3} + \dots)  + \frac{2^{3/2}}{\theta\pi} f(t) \begin{cases} \sqrt{2} & n \geq 1 \\ 1 & n = 0 \end{cases}
  ,
\end{align*}
where we take $f(t - \theta) = 0$ as we no longer have access to it (this holds
when $t < \theta$ as well).

In the usual way, we can write this as linear dynamics
\begin{align*}
  \ddt c(t) &= -\frac{1}{\theta} A c(t) + \frac{1}{\theta} B f(t)
  \\
  A &=
  4
  \begin{bmatrix}
    0 & & & & \dots \\
    2^{-\frac{1}{2}} & 0 & & & \\
    0 & 2 & 0 & & \dots \\
    2^{-\frac{1}{2}} \cdot 3 & 0 & 3 & 0 & \\
    & \ddots & & \ddots &
  \end{bmatrix}
  \\
  B &=
  \frac{2^{3/2}}{\pi}
  \begin{bmatrix}
    1 \\ \sqrt{2} \\ \sqrt{2} \\ \sqrt{2} \\ \vdots
  \end{bmatrix}
\end{align*}



\paragraph{Reconstruction}
In the interval $(t-\theta, t)$,
\begin{align*}
  f(x) \approx \sum_{n=0}^{N-1} c_n(t) p_n(t, x) \chi(t, x).
\end{align*}





\end{document}

